% !TeX spellcheck = es
\documentclass[12pt]{extarticle}
\usepackage[]{geometry}
\newgeometry{left=2cm, right=2cm, top=2cm, bottom=2cm}
\usepackage[parfill]{parskip}
\usepackage[utf8]{inputenc}
\usepackage[24hr]{datetime}
\usepackage{hyperref}
\usepackage{xcolor}
\hypersetup{
    colorlinks,
    linkcolor={red!50!black},
    citecolor={blue!50!black},
    urlcolor={blue!80!black}
}
\usepackage{authoraftertitle}

\title{Línea de trabajo sobre Serverless\\ {\large Credit-based serverless containers resource management}}
\author{Jonatan Enes}
\date{\today~\currenttime}

% ----------------------------------------------------------------------------------%
% --------- ^ This is the common latex code for any doc, keep it simple ^ --------- %
% ----------------------------------------------------------------------------------%


\newcommand{\DepTareaComun}{\textcolor{blue}{Se necesita la Tarea común $\Rightarrow$~}}


% ---------------------------------------------------------------------------------------%
% --------- ^ This is the customized latex code for this doc, keep it simple ^ --------- %
% ---------------------------------------------------------------------------------------%


\begin{document}
    \maketitle

    \tableofcontents

    \section{Descripción}

    Se implementaría un escenario donde el framework aplicaría un control de costes mediante políticas de contabilidad y presupuestado.
    Para el presupuestado se usaría un sistema de crédito, el cual se podría asociar a un contenedor o aplicación.

    \begin{itemize}
        \item Según se utilicen recursos, se irá restando de este crédito.
        \item Cuando se consuma el crédito, se quitan los recursos o se reducen al mínimo.
        \item Si aparece nuevo crédito
        \begin{itemize}
            \item durante la ejecución, se mantienen los recursos
            \item una vez finalizado el crédito, se permitiría consumir recursos de nuevo
        \end{itemize}
    \end{itemize}

    La gestión de este crédito se podría hacer internamente, o conectar con una \textit{`wallet'} de una blockchain.

    \section{Objetivos}
    \begin{itemize}
        \item Ampliar el framework con una política de contabilidad mediante la cual se midan los gastos de un contenedor/aplicación y se anoten
        \item Ampliar el framework con un control de presupuesto, que combinado con la contabilidad, permita asignar balances a un contenedor/aplicación y saber en todo momento los recursos que puede tener
        \item Integrar la parte de presupuestado, más concretamente los balances del crédito, con wallets de una blockchain
    \end{itemize}

    \section{Novedad}

    Aparte de la propia plataforma de contenedores con paradigma serverless, donde los usuarios ejecutan sus aplicaciones sin preocuparse de definir unos recursos iniciales, además se puede establecer un control del consumo de dichas aplicaciones mediante el flujo de entrada de 'crédito'.
    Si se usa una blockchain, se añadiría además la anonimidad. El uso de la blockchain además protege de la deial of wallet, ya que los fondos vienen de fuera.

    \section{Estado del arte}
    \begin{itemize}
    	\item Serverless computing: a security perspective
    	\begin{itemize}
    		\item Describe muchos aspectos de seguridad de los contenedores, algunos pueden ser de utilidad como el denial of wallet
    	\end{itemize}
    	\item Towards Serverless as Commodity: a case of Knative
    	\item Denial of wallet—Defining a looming threat to serverless computing
    \end{itemize}
    
    \section{Experimentos}
    \begin{itemize}
        \item Se arranca un contenedor o aplicación con fondos, se espera hasta que se agoten y se observa como los recursos se restringen
        \item Con un contenedor o aplicación que tiene los fondos agotados, se le da nuevos fondos y se observa como los recursos vuelven a estar disponibles
        \item Si se usa una blockchain, arrancar un contenedor o aplicación sin fondos, hacer una transferencia y observar como se le dan recursos, esperar hasta que se agoten, y darle de nuevo
    \end{itemize}

    \section{Tareas}

    \begin{itemize}
        \item Implementar un servicio o base de datos de algún tipo (blockchain), que guarde los balances de los contenedores/aplicaciones/usuarios
        \item Implementar un servicio que calcule, cada cierto tiempo, cuanto ha gastado un contenedor/aplicación/usuario y lo reste de su balance
        \item Implementar un servicio que compruebe que el balance es positivo, y que sino, restrinja los recursos, o que lleve a cabo las acciones necesarias para esto
    \end{itemize}


\end{document}